% GENERATED FILE. Edit canonical lessons, then run npm run generate:books.
% canonical-source-hash: fnv1a64:578e6b5a8f39e894
% canonical-lessons: TA-C240-urutiyana, TA-C240-vattamana, TA-C240-caturamana, TA-C240-kurmaiyana, TA-C240-arumaiyana

\chapter{Firm, Round, Square, Sharp, Wonderful}
\label{ch:ta-firm-round-square-sharp-wonderful}
\hlchaptermodality{240}

\begin{chapteropening}
\textbf{By the end of this chapter:} \emph{I can say firm, round, square, sharp and wonderful.}

Complete the last lesson of chapter 240: I can say firm, round, square, sharp and wonderful.
\end{chapteropening}

\section[\texorpdfstring{\ta{உறுதியான} (uṟutiyāṉa)}{உறுதியான (uṟutiyāṉa)}]{\ta{உறுதியான} (uṟutiyāṉa) — firm, sturdy}

\label{lesson:TA-C240-urutiyana}



\begin{quote}
Before the new one: say the Tamil for pity, then the Tamil for loneliness.
\end{quote}

\subsection*{You'll want to know: \ta{உறுதியான}}
\textbf{\ta{உறுதியான}} — \emph{uṟutiyāṉa} — \textquotedblleft{}firm, sturdy\textquotedblright{}.

\begin{cousinweb}
\textbf{\ta{உறுதியான}} (\emph{uṟutiyāṉa}) is native Dravidian. DEDR 721 lists Malayalam \emph{uṟekka}, \textquotedblleft{}to be firm\textquotedblright{}. It is \emph{uṟuti}, \textquotedblleft{}firmness\textquotedblright{}, with \emph{-āṉa}.
\end{cousinweb}

\begin{grammarlens}[title={What you've built}]
One more word for this chapter.
\end{grammarlens}

\subsection*{Your turn}
\begin{itemize}
\raggedright
  \item \emph{Say it:} \emph{uṟutiyāṉa}
  \item \emph{Say it:} \emph{uṟutiyāṉa}, once more
  \item \emph{Say it:} say \emph{uṟutiyāṉa}
  \item \emph{Recall:} say the Tamil for pity, then the Tamil for loneliness, then say \emph{uṟutiyāṉa} again
\end{itemize}

\begin{tcolorbox}[breakable,colback=teal!4,colframe=teal!35!black,title={Before you move on}]
What does \ta{உறுதியான} mean? (\textquotedblleft{}Firm, sturdy\textquotedblright{}.) Say it once more.
\end{tcolorbox}
\section[\texorpdfstring{\ta{வட்டமான} (vaṭṭamāṉa)}{வட்டமான (vaṭṭamāṉa)}]{\ta{வட்டமான} (vaṭṭamāṉa) — round}

\label{lesson:TA-C240-vattamana}



\begin{quote}
Before the new one: say the Tamil for loneliness, then the Tamil for firm.
\end{quote}

\subsection*{You'll want to know: \ta{வட்டமான}}
\textbf{\ta{வட்டமான}} — \emph{vaṭṭamāṉa} — \textquotedblleft{}round\textquotedblright{}.

\begin{cousinweb}
\textbf{\ta{வட்டமான}} (\emph{vaṭṭamāṉa}) is \emph{vaṭṭam}, \textquotedblleft{}circle\textquotedblright{}, usually traced to Sanskrit \emph{vṛtta}, \textquotedblleft{}rounded\textquotedblright{} (CDIAL 12069), with \emph{-āṉa}.
\end{cousinweb}

\begin{grammarlens}[title={What you've built}]
One more word for this chapter.
\end{grammarlens}

\subsection*{Your turn}
\begin{itemize}
\raggedright
  \item \emph{Say it:} \emph{vaṭṭamāṉa}
  \item \emph{Say it:} \emph{vaṭṭamāṉa}, once more
  \item \emph{Say it:} say \emph{vaṭṭamāṉa}
  \item \emph{Recall:} say the Tamil for loneliness, then the Tamil for firm, then say \emph{vaṭṭamāṉa} again
\end{itemize}

\begin{tcolorbox}[breakable,colback=teal!4,colframe=teal!35!black,title={Before you move on}]
What does \ta{வட்டமான} mean? (\textquotedblleft{}Round\textquotedblright{}.) Say it once more.
\end{tcolorbox}
\section[\texorpdfstring{\ta{சதுரமான} (caturamāṉa)}{சதுரமான (caturamāṉa)}]{\ta{சதுரமான} (caturamāṉa) — square}

\label{lesson:TA-C240-caturamana}



\begin{quote}
Before the new one: say the Tamil for firm, then the Tamil for round.
\end{quote}

\subsection*{You'll want to know: \ta{சதுரமான}}
\textbf{\ta{சதுரமான}} — \emph{caturamāṉa} — \textquotedblleft{}square\textquotedblright{}.

\begin{cousinweb}
\textbf{\ta{சதுரமான}} (\emph{caturamāṉa}) is \emph{caturam}, \textquotedblleft{}square\textquotedblright{}, a well-established loan from Sanskrit \emph{caturasra}, \textquotedblleft{}four-cornered\textquotedblright{}, with \emph{-āṉa}.
\end{cousinweb}

\begin{grammarlens}[title={What you've built}]
One more word for this chapter.
\end{grammarlens}

\subsection*{Your turn}
\begin{itemize}
\raggedright
  \item \emph{Say it:} \emph{caturamāṉa}
  \item \emph{Say it:} \emph{caturamāṉa}, once more
  \item \emph{Say it:} say \emph{caturamāṉa}
  \item \emph{Recall:} say the Tamil for firm, then the Tamil for round, then say \emph{caturamāṉa} again
\end{itemize}

\begin{tcolorbox}[breakable,colback=teal!4,colframe=teal!35!black,title={Before you move on}]
What does \ta{சதுரமான} mean? (\textquotedblleft{}Square\textquotedblright{}.) Say it once more.
\end{tcolorbox}
\section[\texorpdfstring{\ta{கூர்மையான} (kūrmaiyāṉa)}{கூர்மையான (kūrmaiyāṉa)}]{\ta{கூர்மையான} (kūrmaiyāṉa) — sharp}

\label{lesson:TA-C240-kurmaiyana}



\begin{quote}
Before the new one: say the Tamil for round, then the Tamil for square.
\end{quote}

\subsection*{You'll want to know: \ta{கூர்மையான}}
\textbf{\ta{கூர்மையான}} — \emph{kūrmaiyāṉa} — \textquotedblleft{}sharp\textquotedblright{}.

\begin{cousinweb}
\textbf{\ta{கூர்மையான}} (\emph{kūrmaiyāṉa}) is native Dravidian. DEDR 1898 lists Malayalam \emph{kūr}, \textquotedblleft{}sharpness\textquotedblright{}.
\end{cousinweb}

\begin{grammarlens}[title={What you've built}]
One more word for this chapter.
\end{grammarlens}

\subsection*{Your turn}
\begin{itemize}
\raggedright
  \item \emph{Say it:} \emph{kūrmaiyāṉa}
  \item \emph{Say it:} \emph{kūrmaiyāṉa}, once more
  \item \emph{Say it:} say \emph{kūrmaiyāṉa}
  \item \emph{Recall:} say the Tamil for round, then the Tamil for square, then say \emph{kūrmaiyāṉa} again
\end{itemize}

\begin{tcolorbox}[breakable,colback=teal!4,colframe=teal!35!black,title={Before you move on}]
What does \ta{கூர்மையான} mean? (\textquotedblleft{}Sharp\textquotedblright{}.) Say it once more.
\end{tcolorbox}
\section[\texorpdfstring{\ta{அருமையான} (arumaiyāṉa)}{அருமையான (arumaiyāṉa)}]{\ta{அருமையான} (arumaiyāṉa) — wonderful, excellent}

\label{lesson:TA-C240-arumaiyana}



\begin{quote}
Before the new one: say the Tamil for square, then the Tamil for sharp.
\end{quote}

\subsection*{You'll want to know: \ta{அருமையான}}
\textbf{\ta{அருமையான}} — \emph{arumaiyāṉa} — \textquotedblleft{}wonderful, excellent\textquotedblright{}.

\begin{cousinweb}
\textbf{\ta{அருமையான}} (\emph{arumaiyāṉa}) is native Dravidian. DEDR 221 lists Malayalam \emph{aruma}, \textquotedblleft{}rarity\textquotedblright{}. It is \emph{arumai}, \textquotedblleft{}rareness, excellence\textquotedblright{}, with \emph{-āṉa}: what is rare is precious.
\end{cousinweb}

\begin{grammarlens}[title={What you've built}]
One more word for this chapter.
\end{grammarlens}

\subsection*{Your turn}
\begin{itemize}
\raggedright
  \item \emph{Say it:} \emph{arumaiyāṉa}
  \item \emph{Say it:} \emph{arumaiyāṉa}, once more
  \item \emph{Say it:} say \emph{arumaiyāṉa}
  \item \emph{Recall:} say the Tamil for square, then the Tamil for sharp, then say \emph{arumaiyāṉa} again
\end{itemize}

\begin{tcolorbox}[breakable,colback=teal!4,colframe=teal!35!black,title={Before you move on}]
What does \ta{அருமையான} mean? (\textquotedblleft{}Wonderful, excellent\textquotedblright{}.) Say it once more.
\end{tcolorbox}
